% ---------------------------------------------------------
% TABELA EMENTA 14: Física — Ano 1
% ---------------------------------------------------------
\begin{xltabular}{\linewidth}{|X|p{2.5cm}|p{2.5cm}|}
\hline
\multirow{3}{=}{\textcolor{ifscgreen}{\textbf{Unidade Curricular:}}\\[3pt]\textbf{\large Física — Ano 1}} & \multicolumn{2}{p{\dimexpr5cm+2\tabcolsep+\arrayrulewidth\relax}|}{\textcolor{ifscgreen}{\textbf{Semestre:}} \textbf{2º (Física I)}} \\ \cline{2-3}
 & \textcolor{ifscgreen}{\textbf{CH EaD*:}} & \textcolor{ifscgreen}{\textbf{CH Total*:}} \\ \cline{2-3}
 & \textbf{00 h} & \textbf{40 h} \\ \hline
\endhead
\multicolumn{3}{|p{\dimexpr\linewidth-2\tabcolsep-2\arrayrulewidth\relax}|}{\textcolor{ifscgreen}{\textbf{Objetivos:}}} \\ \hline
\multicolumn{3}{|p{\dimexpr\linewidth-2\tabcolsep-2\arrayrulewidth\relax}|}{
\begin{itemize}[leftmargin=*,noitemsep,topsep=2pt]
 \item Compreender a Física como ciência, reconhecendo sua evolução histórica, seus principais ramos e sua importância para a compreensão dos fenômenos naturais e para o desenvolvimento científico e tecnológico.
 \item Identificar e utilizar corretamente as grandezas físicas fundamentais e derivadas, bem como o Sistema Internacional de Unidades (SI) e suas conversões.
 \item Desenvolver a capacidade de interpretar fenômenos físicos por meio da linguagem matemática, da análise de gráficos, tabelas e relações entre grandezas.
 \item Compreender os conceitos fundamentais da Cinemática, descrevendo e analisando o movimento de corpos em diferentes referenciais.
 \item Aplicar as equações do Movimento Uniforme (MU) e do Movimento Uniformemente Variado (MUV) na resolução de problemas envolvendo deslocamento, velocidade e aceleração.
 \item Interpretar e construir gráficos de posição, velocidade e aceleração em função do tempo.
 \item Compreender os conceitos fundamentais da Dinâmica, identificando as forças que atuam sobre um corpo e seus efeitos sobre o movimento.
 \item Analisar e aplicar as Leis de Newton na interpretação e resolução de problemas relacionados ao equilíbrio e ao movimento de corpos.
 \item Desenvolver o raciocínio lógico e quantitativo na resolução de problemas físicos.
 \item Relacionar os conceitos estudados com situações do cotidiano, reconhecendo os princípios da Mecânica em diferentes contextos científicos, tecnológicos e sociais.
\end{itemize}
} \\ \hline
\multicolumn{3}{|p{\dimexpr\linewidth-2\tabcolsep-2\arrayrulewidth\relax}|}{\textcolor{ifscgreen}{\textbf{Conteúdos:}}} \\ \hline
\multicolumn{3}{|p{\dimexpr\linewidth-2\tabcolsep-2\arrayrulewidth\relax}|}{
\begin{itemize}[leftmargin=*,noitemsep,topsep=2pt]
 \item Introdução à Física: a origem da Física, as áreas da Física, grandezas físicas e unidades de medidas; noções básicas de Cinemática; Movimento Uniforme e Movimento Uniformemente Variado; Dinâmica (introdução); Leis de Newton e suas aplicações.
\end{itemize}
} \\ \hline
\multicolumn{3}{|p{\dimexpr\linewidth-2\tabcolsep-2\arrayrulewidth\relax}|}{\textcolor{ifscgreen}{\textbf{Estratégias de Ensino e Aprendizagem:}}} \\ \hline
\multicolumn{3}{|p{\dimexpr\linewidth-2\tabcolsep-2\arrayrulewidth\relax}|}{- **Metodologia:** Aulas majoritariamente expositivas e dialogadas, sob proposta dialógico-problematizadora a partir de fenômenos naturais, tecnológicos e produtivos do mundo do trabalho, com três momentos pedagógicos: problematização inicial, organização e aplicação dos conhecimentos. Conteúdos contextualizados e interdisciplinares, com resolução de listas de exercícios; exposição de vídeos; seminários; trabalhos de pesquisa; montagem e apresentação de experimentos; desenvolvimento de projetos; interpretação de textos técnicos e científicos. Recursos: livros didáticos e digitais, apostilas, projetor multimídia, equipamentos de laboratório.
- **Avaliação:** Perspectiva formativa, levando em conta todas as atividades realizadas pelos estudantes.} \\ \hline
\multicolumn{3}{|p{\dimexpr\linewidth-2\tabcolsep-2\arrayrulewidth\relax}|}{\textcolor{ifscgreen}{\textbf{Bibliografia Básica:}}} \\ \hline
\multicolumn{3}{|p{\dimexpr\linewidth-2\tabcolsep-2\arrayrulewidth\relax}|}{1. Livro didático fornecido pelo Fundo Nacional de Desenvolvimento da Educação (FNDE).
2. HEWITT, Paul G. \textbf{Física conceitual}. 12. ed. Porto Alegre: Bookman, 2015.
3. LUZ, Antônio Máximo Ribeiro da; ALVARENGA, Beatriz Gonçalves de. \textbf{Curso de física, volume 1}. 6. ed. rev. ampl. São Paulo: Scipione, 2006.} \\ \hline
\multicolumn{3}{|p{\dimexpr\linewidth-2\tabcolsep-2\arrayrulewidth\relax}|}{\textcolor{ifscgreen}{\textbf{Bibliografia Complementar:}}} \\ \hline
\multicolumn{3}{|p{\dimexpr\linewidth-2\tabcolsep-2\arrayrulewidth\relax}|}{1. BARRETO FILHO, Benigno; SILVA, Cláudio Xavier da. \textbf{360º: física: aula por aula: volume único}. 3. ed. São Paulo: FTD, 2015.
2. YAMAMOTO, Kazuhito; FUKE, Luiz Felipe. \textbf{Física para o ensino médio 1: mecânica}. 3. ed. São Paulo: Saraiva, 2013.
3. KNIGHT, Randall D. \textbf{Física: uma abordagem estratégica: volume 1: mecânica newtoniana, gravitação, oscilações e ondas}. 2. ed. Porto Alegre: Bookman, 2009.} \\ \hline
\end{xltabular}

\vspace{0.4cm}


% ---------------------------------------------------------
% TABELA EMENTA 15: Física — Ano 2
% ---------------------------------------------------------
\begin{xltabular}{\linewidth}{|X|p{2.5cm}|p{2.5cm}|}
\hline
\multirow{3}{=}{\textcolor{ifscgreen}{\textbf{Unidade Curricular:}}\\[3pt]\textbf{\large Física — Ano 2}} & \multicolumn{2}{p{\dimexpr5cm+2\tabcolsep+\arrayrulewidth\relax}|}{\textcolor{ifscgreen}{\textbf{Semestre:}} \textbf{3º e 4º (Física II)}} \\ \cline{2-3}
 & \textcolor{ifscgreen}{\textbf{CH EaD*:}} & \textcolor{ifscgreen}{\textbf{CH Total*:}} \\ \cline{2-3}
 & \textbf{00 h} & \textbf{80 h} \\ \hline
\endhead
\multicolumn{3}{|p{\dimexpr\linewidth-2\tabcolsep-2\arrayrulewidth\relax}|}{\textcolor{ifscgreen}{\textbf{Objetivos:}}} \\ \hline
\multicolumn{3}{|p{\dimexpr\linewidth-2\tabcolsep-2\arrayrulewidth\relax}|}{
\begin{itemize}[leftmargin=*,noitemsep,topsep=2pt]
 \item Compreender os conceitos de trabalho, energia e potência, analisando as transformações de energia e aplicando o princípio da conservação da energia mecânica.
 \item Analisar os conceitos de impulso e quantidade de movimento, aplicando o princípio da conservação da quantidade de movimento em colisões e interações entre corpos.
 \item Compreender os fundamentos da Gravitação Universal, relacionando a interação gravitacional aos movimentos de corpos celestes e a fenômenos do cotidiano.
 \item Identificar e aplicar os princípios da Hidrostática, analisando pressão, empuxo e equilíbrio dos fluidos em repouso.
 \item Compreender os conceitos de temperatura, calor, dilatação térmica e comportamento dos gases ideais.
 \item Aplicar as Leis da Termodinâmica na análise das transformações de energia e das máquinas térmicas.
 \item Compreender os fenômenos da reflexão e da refração da luz, analisando a formação de imagens.
 \item Compreender os princípios do movimento ondulatório, relacionando reflexão, refração, difração e interferência a situações do cotidiano.
 \item Desenvolver o raciocínio lógico, quantitativo e científico por meio da resolução de problemas.
 \item Relacionar os conceitos de Mecânica, Termologia, Óptica e Ondulatória com fenômenos naturais e aplicações tecnológicas.
\end{itemize}
} \\ \hline
\multicolumn{3}{|p{\dimexpr\linewidth-2\tabcolsep-2\arrayrulewidth\relax}|}{\textcolor{ifscgreen}{\textbf{Conteúdos:}}} \\ \hline
\multicolumn{3}{|p{\dimexpr\linewidth-2\tabcolsep-2\arrayrulewidth\relax}|}{
\begin{itemize}[leftmargin=*,noitemsep,topsep=2pt]
 \item Trabalho de uma força e Conservação da Energia Mecânica; Impulso e Conservação da quantidade de movimento; Gravitação Universal; Hidrostática; Termometria, Dilatação Térmica e Gás Ideal; Leis da Termodinâmica e Máquinas Térmicas; Reflexão da Luz e Refração da Luz; Movimento Ondulatório.
\end{itemize}
} \\ \hline
\multicolumn{3}{|p{\dimexpr\linewidth-2\tabcolsep-2\arrayrulewidth\relax}|}{\textcolor{ifscgreen}{\textbf{Estratégias de Ensino e Aprendizagem:}}} \\ \hline
\multicolumn{3}{|p{\dimexpr\linewidth-2\tabcolsep-2\arrayrulewidth\relax}|}{- **Metodologia:** Aulas majoritariamente expositivas e dialogadas, sob proposta dialógico-problematizadora a partir de fenômenos naturais, tecnológicos e produtivos do mundo do trabalho, com três momentos pedagógicos: problematização inicial, organização e aplicação dos conhecimentos. Conteúdos contextualizados e interdisciplinares, com resolução de listas de exercícios; exposição de vídeos; seminários; trabalhos de pesquisa; montagem e apresentação de experimentos; desenvolvimento de projetos; interpretação de textos técnicos e científicos.
- **Avaliação:** Perspectiva formativa, levando em conta todas as atividades realizadas pelos estudantes.} \\ \hline
\multicolumn{3}{|p{\dimexpr\linewidth-2\tabcolsep-2\arrayrulewidth\relax}|}{\textcolor{ifscgreen}{\textbf{Bibliografia Básica:}}} \\ \hline
\multicolumn{3}{|p{\dimexpr\linewidth-2\tabcolsep-2\arrayrulewidth\relax}|}{1. Livro didático fornecido pelo Fundo Nacional de Desenvolvimento da Educação (FNDE).
2. HEWITT, Paul G. \textbf{Física conceitual}. 12. ed. Porto Alegre: Bookman, 2015.
3. LUZ, Antônio Máximo Ribeiro da; ALVARENGA, Beatriz Gonçalves de. \textbf{Curso de física, volume 1}. 6. ed. rev. ampl. São Paulo: Scipione, 2006.
4. LUZ, Antônio Máximo Ribeiro da; ALVARENGA, Beatriz Gonçalves de. \textbf{Curso de física, volume 2}. 6. ed. rev. ampl. São Paulo: Scipione, 2006.} \\ \hline
\multicolumn{3}{|p{\dimexpr\linewidth-2\tabcolsep-2\arrayrulewidth\relax}|}{\textcolor{ifscgreen}{\textbf{Bibliografia Complementar:}}} \\ \hline
\multicolumn{3}{|p{\dimexpr\linewidth-2\tabcolsep-2\arrayrulewidth\relax}|}{1. BARRETO FILHO, Benigno; SILVA, Cláudio Xavier da. \textbf{360º: física: aula por aula: volume único}. 3. ed. São Paulo: FTD, 2015.
2. YAMAMOTO, Kazuhito; FUKE, Luiz Felipe. \textbf{Física para o ensino médio 2: termologia, óptica, ondulatória}. 3. ed. São Paulo: Saraiva, 2013.
3. KNIGHT, Randall D. \textbf{Física: uma abordagem estratégica: volume 2: termodinâmica e óptica}. 2. ed. Porto Alegre: Bookman, 2009.} \\ \hline
\end{xltabular}

\vspace{0.4cm}


% ---------------------------------------------------------
% TABELA EMENTA 16: Física — Ano 3
% ---------------------------------------------------------
\begin{xltabular}{\linewidth}{|X|p{2.5cm}|p{2.5cm}|}
\hline
\multirow{3}{=}{\textcolor{ifscgreen}{\textbf{Unidade Curricular:}}\\[3pt]\textbf{\large Física — Ano 3}} & \multicolumn{2}{p{\dimexpr5cm+2\tabcolsep+\arrayrulewidth\relax}|}{\textcolor{ifscgreen}{\textbf{Semestre:}} \textbf{4º e 5º (Física III)}} \\ \cline{2-3}
 & \textcolor{ifscgreen}{\textbf{CH EaD*:}} & \textcolor{ifscgreen}{\textbf{CH Total*:}} \\ \cline{2-3}
 & \textbf{00 h} & \textbf{80 h} \\ \hline
\endhead
\multicolumn{3}{|p{\dimexpr\linewidth-2\tabcolsep-2\arrayrulewidth\relax}|}{\textcolor{ifscgreen}{\textbf{Objetivos:}}} \\ \hline
\multicolumn{3}{|p{\dimexpr\linewidth-2\tabcolsep-2\arrayrulewidth\relax}|}{
\begin{itemize}[leftmargin=*,noitemsep,topsep=2pt]
 \item Compreender os conceitos fundamentais da eletrostática, analisando a interação entre cargas elétricas por meio da força eletrostática, do campo elétrico, do potencial elétrico e da energia potencial elétrica.
 \item Aplicar os conceitos de corrente elétrica, tensão, resistência e potência na análise e resolução de problemas envolvendo circuitos elétricos de corrente contínua (CC) e corrente alternada (CA).
 \item Identificar as características e aplicações dos principais componentes de circuitos elétricos, interpretando esquemas elétricos e analisando associações de resistores em série e em paralelo.
 \item Compreender os princípios do magnetismo, analisando a interação entre campos magnéticos, correntes elétricas e materiais magnéticos.
 \item Explicar os fenômenos da indução eletromagnética, relacionando-os ao funcionamento de geradores, transformadores e motores elétricos.
 \item Compreender a natureza e as propriedades das ondas eletromagnéticas e suas aplicações tecnológicas.
 \item Identificar os fundamentos da Física Moderna, reconhecendo as limitações da Física Clássica e os princípios básicos da Relatividade e da Mecânica Quântica.
 \item Desenvolver o raciocínio lógico, quantitativo e científico relacionado aos fenômenos eletromagnéticos e modernos.
 \item Relacionar os conceitos de Eletrostática, Eletrodinâmica, Magnetismo, Ondas Eletromagnéticas e Física Moderna com fenômenos naturais e aplicações tecnológicas.
\end{itemize}
} \\ \hline
\multicolumn{3}{|p{\dimexpr\linewidth-2\tabcolsep-2\arrayrulewidth\relax}|}{\textcolor{ifscgreen}{\textbf{Conteúdos:}}} \\ \hline
\multicolumn{3}{|p{\dimexpr\linewidth-2\tabcolsep-2\arrayrulewidth\relax}|}{
\begin{itemize}[leftmargin=*,noitemsep,topsep=2pt]
 \item Carga elétrica, Força Eletrostática, Campo Elétrico, Potencial Elétrico e Energia Potencial Elétrica; Corrente elétrica, Tipos de Corrente Elétrica, Resistores e Circuitos Elétricos CC e CA; Diagramas de circuitos — unifilar e multifilar (aplicações); Magnetismo e Indução Eletromagnética; Ondas Eletromagnéticas; Introdução à Física Moderna.
\end{itemize}
} \\ \hline
\multicolumn{3}{|p{\dimexpr\linewidth-2\tabcolsep-2\arrayrulewidth\relax}|}{\textcolor{ifscgreen}{\textbf{Estratégias de Ensino e Aprendizagem:}}} \\ \hline
\multicolumn{3}{|p{\dimexpr\linewidth-2\tabcolsep-2\arrayrulewidth\relax}|}{- **Metodologia:** Aulas majoritariamente expositivas e dialogadas, sob proposta dialógico-problematizadora a partir de fenômenos naturais, tecnológicos e produtivos do mundo do trabalho, com três momentos pedagógicos: problematização inicial, organização e aplicação dos conhecimentos. Exposição de vídeos; seminários; trabalhos de pesquisa; montagem e apresentação de experimentos; desenvolvimento de projetos; interpretação de textos técnicos e científicos.
- **Avaliação:** Perspectiva formativa, levando em conta todas as atividades realizadas pelos estudantes.} \\ \hline
\multicolumn{3}{|p{\dimexpr\linewidth-2\tabcolsep-2\arrayrulewidth\relax}|}{\textcolor{ifscgreen}{\textbf{Bibliografia Básica:}}} \\ \hline
\multicolumn{3}{|p{\dimexpr\linewidth-2\tabcolsep-2\arrayrulewidth\relax}|}{1. Livro didático fornecido pelo Fundo Nacional de Desenvolvimento da Educação (FNDE).
2. HEWITT, Paul G. \textbf{Física conceitual}. 12. ed. Porto Alegre: Bookman, 2015.
3. LUZ, Antônio Máximo Ribeiro da; ALVARENGA, Beatriz Gonçalves de. \textbf{Curso de física, volume 3}. 6. ed. rev. ampl. São Paulo: Scipione, 2006.} \\ \hline
\multicolumn{3}{|p{\dimexpr\linewidth-2\tabcolsep-2\arrayrulewidth\relax}|}{\textcolor{ifscgreen}{\textbf{Bibliografia Complementar:}}} \\ \hline
\multicolumn{3}{|p{\dimexpr\linewidth-2\tabcolsep-2\arrayrulewidth\relax}|}{1. BARRETO FILHO, Benigno; SILVA, Cláudio Xavier da. \textbf{360º: física: aula por aula: volume único}. 3. ed. São Paulo: FTD, 2015.
2. YAMAMOTO, Kazuhito; FUKE, Luiz Felipe. \textbf{Física para o ensino médio 3: eletricidade, física moderna}. 3. ed. São Paulo: Saraiva, 2013.
3. KNIGHT, Randall D. \textbf{Física: uma abordagem estratégica: volume 3: eletricidade e magnetismo}. 2. ed. Porto Alegre: Bookman, 2009.} \\ \hline
\end{xltabular}

\vspace{0.4cm}


% ---------------------------------------------------------
% TABELA EMENTA 17: Matemática — Ano 1
% ---------------------------------------------------------
\begin{xltabular}{\linewidth}{|X|p{2.5cm}|p{2.5cm}|}
\hline
\multirow{3}{=}{\textcolor{ifscgreen}{\textbf{Unidade Curricular:}}\\[3pt]\textbf{\large Matemática — Ano 1}} & \multicolumn{2}{p{\dimexpr5cm+2\tabcolsep+\arrayrulewidth\relax}|}{\textcolor{ifscgreen}{\textbf{Semestre:}} \textbf{1º e 2º}} \\ \cline{2-3}
 & \textcolor{ifscgreen}{\textbf{CH EaD*:}} & \textcolor{ifscgreen}{\textbf{CH Total*:}} \\ \cline{2-3}
 & \textbf{00 h} & \textbf{80 h} \\ \hline
\endhead
\multicolumn{3}{|p{\dimexpr\linewidth-2\tabcolsep-2\arrayrulewidth\relax}|}{\textcolor{ifscgreen}{\textbf{Objetivos:}}} \\ \hline
\multicolumn{3}{|p{\dimexpr\linewidth-2\tabcolsep-2\arrayrulewidth\relax}|}{
\begin{itemize}[leftmargin=*,noitemsep,topsep=2pt]
 \item Compreender e utilizar adequadamente a linguagem matemática na resolução de problemas, relacionando-a ao contexto da área de Administração.
 \item Analisar, interpretar e utilizar os conhecimentos elencados pela disciplina na resolução de problemas relacionados à área de Administração.
\end{itemize}
} \\ \hline
\multicolumn{3}{|p{\dimexpr\linewidth-2\tabcolsep-2\arrayrulewidth\relax}|}{\textcolor{ifscgreen}{\textbf{Conteúdos:}}} \\ \hline
\multicolumn{3}{|p{\dimexpr\linewidth-2\tabcolsep-2\arrayrulewidth\relax}|}{
\begin{itemize}[leftmargin=*,noitemsep,topsep=2pt]
 \item \textbf{Parte I — Fundamentos:} Potenciação (definição, propriedades, expoentes fracionários); Radiciação (definição, propriedades, simplificação, racionalização, equações com radicais, gráficos e aplicações); Conjuntos numéricos (naturais, inteiros, racionais, irracionais e reais, operações e propriedades, conversão entre decimais e frações, dízimas periódicas e não periódicas); Equações do 1º grau; Equações do 2º grau; Sistemas lineares 2x2 (método da adição e substituição).
 \item \textbf{Parte II — Matemática Ensino Médio:} Conjuntos (operações, diagrama de Venn); Funções (definição, crescimento/decrescimento, domínio, contradomínio e imagem); Função afim (gráfico, coeficientes, aplicações); Função quadrática (gráfico, vértice, raízes, concavidade, aplicações); Teorema de Pitágoras e razões trigonométricas; Noções de estatística (coleta de dados, tabelas e gráficos, medidas de tendência central e desvio padrão).
\end{itemize}
} \\ \hline
\multicolumn{3}{|p{\dimexpr\linewidth-2\tabcolsep-2\arrayrulewidth\relax}|}{\textcolor{ifscgreen}{\textbf{Estratégias de Ensino e Aprendizagem:}}} \\ \hline
\multicolumn{3}{|p{\dimexpr\linewidth-2\tabcolsep-2\arrayrulewidth\relax}|}{- **Metodologia:** Aulas expositivas dialogadas; aulas de exercícios; apresentação em linguagem verbal por escrito ou em diálogos na sala de aula; discussões em grupos; estudos dirigidos; leitura e interpretação por meio de datashow; pesquisas em laboratório de informática; seminários; trabalho em laboratório de informática e na biblioteca; trabalhos individuais e em grupos; uso de jogos e objetos de aprendizagem; experimentação no ensino de matemática; produção de vídeos de matemática por estudantes.
- **Avaliação:** Avaliações qualitativas e quantitativas durante o semestre.} \\ \hline
\multicolumn{3}{|p{\dimexpr\linewidth-2\tabcolsep-2\arrayrulewidth\relax}|}{\textcolor{ifscgreen}{\textbf{Bibliografia Básica:}}} \\ \hline
\multicolumn{3}{|p{\dimexpr\linewidth-2\tabcolsep-2\arrayrulewidth\relax}|}{1. BONJORNO, José Roberto; GIOVANNI JÚNIOR, José Ruy; GIOVANNI, José Ruy. \textbf{Matemática fundamental: uma nova abordagem: ensino médio: 1º ano}. 3. ed. São Paulo: FTD, 2013.
2. DANTE, Luiz Roberto. \textbf{Projeto Múltiplo: matemática: ensino médio, 1º ano}. São Paulo: Ática, 2014.
3. GIOVANNI, José Ruy et al. \textbf{360º matemática fundamental: uma nova abordagem: volume único}. 2. ed. São Paulo: FTD, 2015.} \\ \hline
\multicolumn{3}{|p{\dimexpr\linewidth-2\tabcolsep-2\arrayrulewidth\relax}|}{\textcolor{ifscgreen}{\textbf{Bibliografia Complementar:}}} \\ \hline
\multicolumn{3}{|p{\dimexpr\linewidth-2\tabcolsep-2\arrayrulewidth\relax}|}{1. LIMA, E. L.; CARVALHO, Paulo Cesar P. \textbf{A Matemática do ensino médio: volume único}. 7. ed. Rio de Janeiro: SBM, 2016.
2. PAIVA, Manoel. \textbf{Matemática: Paiva: 1º ano}. 2. ed. São Paulo: Moderna, 2013.} \\ \hline
\end{xltabular}

\vspace{0.4cm}


% ---------------------------------------------------------
% TABELA EMENTA 18: Matemática — Ano 2
% ---------------------------------------------------------
\begin{xltabular}{\linewidth}{|X|p{2.5cm}|p{2.5cm}|}
\hline
\multirow{3}{=}{\textcolor{ifscgreen}{\textbf{Unidade Curricular:}}\\[3pt]\textbf{\large Matemática — Ano 2}} & \multicolumn{2}{p{\dimexpr5cm+2\tabcolsep+\arrayrulewidth\relax}|}{\textcolor{ifscgreen}{\textbf{Semestre:}} \textbf{3º e 4º}} \\ \cline{2-3}
 & \textcolor{ifscgreen}{\textbf{CH EaD*:}} & \textcolor{ifscgreen}{\textbf{CH Total*:}} \\ \cline{2-3}
 & \textbf{00 h} & \textbf{80 h} \\ \hline
\endhead
\multicolumn{3}{|p{\dimexpr\linewidth-2\tabcolsep-2\arrayrulewidth\relax}|}{\textcolor{ifscgreen}{\textbf{Objetivos:}}} \\ \hline
\multicolumn{3}{|p{\dimexpr\linewidth-2\tabcolsep-2\arrayrulewidth\relax}|}{
\begin{itemize}[leftmargin=*,noitemsep,topsep=2pt]
 \item Compreender e utilizar adequadamente a linguagem matemática na resolução de problemas, relacionando-a ao contexto da área de Administração.
 \item Analisar, interpretar e utilizar os conhecimentos elencados pela disciplina na resolução de problemas relacionados à área de Administração.
\end{itemize}
} \\ \hline
\multicolumn{3}{|p{\dimexpr\linewidth-2\tabcolsep-2\arrayrulewidth\relax}|}{\textcolor{ifscgreen}{\textbf{Conteúdos:}}} \\ \hline
\multicolumn{3}{|p{\dimexpr\linewidth-2\tabcolsep-2\arrayrulewidth\relax}|}{
\begin{itemize}[leftmargin=*,noitemsep,topsep=2pt]
 \item Equação exponencial. Função exponencial (definição, propriedades, gráfico, aplicações) e sua inversa (função logarítmica). Logaritmos (definição e propriedades). Equação logarítmica. Função logarítmica (propriedades, gráfico, aplicações). Círculo trigonométrico (coordenadas, ângulos, quadrantes, seno, cosseno, tangente, periodicidade, aplicações). Funções trigonométricas (seno, cosseno, tangente, secante, cossecante, cotangente, gráficos, aplicações). Identidades trigonométricas. Área e perímetro de figuras planas (triângulos, quadriláteros, hexágono regular e círculo). Geometria espacial (área das superfícies/planificação e volume de prismas, pirâmides, cilindro, cone e esfera).
\end{itemize}
} \\ \hline
\multicolumn{3}{|p{\dimexpr\linewidth-2\tabcolsep-2\arrayrulewidth\relax}|}{\textcolor{ifscgreen}{\textbf{Estratégias de Ensino e Aprendizagem:}}} \\ \hline
\multicolumn{3}{|p{\dimexpr\linewidth-2\tabcolsep-2\arrayrulewidth\relax}|}{- **Metodologia:** Aulas expositivas dialogadas; aulas de exercícios; apresentação em linguagem verbal por escrito ou em diálogos na sala de aula; discussões em grupos; estudos dirigidos; leitura e interpretação por meio de datashow; pesquisas em laboratório de informática; seminários; trabalho em laboratório de informática e na biblioteca; trabalhos individuais e em grupos; uso de jogos e objetos de aprendizagem; experimentação no ensino de matemática; produção de vídeos de matemática por estudantes.
- **Avaliação:** Avaliações qualitativas e quantitativas durante o semestre.} \\ \hline
\multicolumn{3}{|p{\dimexpr\linewidth-2\tabcolsep-2\arrayrulewidth\relax}|}{\textcolor{ifscgreen}{\textbf{Bibliografia Básica:}}} \\ \hline
\multicolumn{3}{|p{\dimexpr\linewidth-2\tabcolsep-2\arrayrulewidth\relax}|}{1. BONJORNO, José Roberto; GIOVANNI JÚNIOR, José Ruy; GIOVANNI, José Ruy. \textbf{Matemática fundamental: uma nova abordagem: ensino médio: volume único}. 3. ed. São Paulo: FTD, 2013.
2. DANTE, Luiz Roberto. \textbf{Projeto Múltiplo: matemática: ensino médio, 2º ano}. São Paulo: Ática, 2014.
3. GIOVANNI, José Ruy et al. \textbf{360º matemática fundamental: uma nova abordagem: volume único}. 2. ed. São Paulo: FTD, 2015.} \\ \hline
\multicolumn{3}{|p{\dimexpr\linewidth-2\tabcolsep-2\arrayrulewidth\relax}|}{\textcolor{ifscgreen}{\textbf{Bibliografia Complementar:}}} \\ \hline
\multicolumn{3}{|p{\dimexpr\linewidth-2\tabcolsep-2\arrayrulewidth\relax}|}{1. LIMA, E. L.; CARVALHO, Paulo Cesar P. \textbf{A Matemática do ensino médio: volume único}. 7. ed. Rio de Janeiro: SBM, 2016.
2. PAIVA, Manoel. \textbf{Matemática: Paiva: 2º ano}. 2. ed. São Paulo: Moderna, 2013.} \\ \hline
\end{xltabular}

\vspace{0.4cm}


% ---------------------------------------------------------
% TABELA EMENTA 19: Matemática — Ano 3
% ---------------------------------------------------------
\begin{xltabular}{\linewidth}{|X|p{2.5cm}|p{2.5cm}|}
\hline
\multirow{3}{=}{\textcolor{ifscgreen}{\textbf{Unidade Curricular:}}\\[3pt]\textbf{\large Matemática — Ano 3}} & \multicolumn{2}{p{\dimexpr5cm+2\tabcolsep+\arrayrulewidth\relax}|}{\textcolor{ifscgreen}{\textbf{Semestre:}} \textbf{5º e 6º}} \\ \cline{2-3}
 & \textcolor{ifscgreen}{\textbf{CH EaD*:}} & \textcolor{ifscgreen}{\textbf{CH Total*:}} \\ \cline{2-3}
 & \textbf{00 h} & \textbf{80 h} \\ \hline
\endhead
\multicolumn{3}{|p{\dimexpr\linewidth-2\tabcolsep-2\arrayrulewidth\relax}|}{\textcolor{ifscgreen}{\textbf{Objetivos:}}} \\ \hline
\multicolumn{3}{|p{\dimexpr\linewidth-2\tabcolsep-2\arrayrulewidth\relax}|}{
\begin{itemize}[leftmargin=*,noitemsep,topsep=2pt]
 \item Compreender e utilizar adequadamente a linguagem matemática na resolução de problemas, relacionando-a ao contexto da área de Administração.
 \item Analisar, interpretar e utilizar os conhecimentos elencados pela disciplina na resolução de problemas relacionados à área de Administração.
\end{itemize}
} \\ \hline
\multicolumn{3}{|p{\dimexpr\linewidth-2\tabcolsep-2\arrayrulewidth\relax}|}{\textcolor{ifscgreen}{\textbf{Conteúdos:}}} \\ \hline
\multicolumn{3}{|p{\dimexpr\linewidth-2\tabcolsep-2\arrayrulewidth\relax}|}{
\begin{itemize}[leftmargin=*,noitemsep,topsep=2pt]
 \item Matrizes (definição, tipos, operações, matriz inversa e aplicações). Determinante de ordens 1, 2 e 3 (solução por Sarrus, propriedades). Sistemas lineares (sistemas equivalentes e homogêneos, resolução por escalonamento, sistemas com variáveis livres). Geometria analítica (distância entre pontos, estudo das retas). Análise combinatória (princípio da contagem, fatorial, arranjos, permutações e combinações). Probabilidade (espaço amostral e evento, probabilidade, união de probabilidade, probabilidade condicional). Polinômios (função polinomial, valor numérico, polinômio nulo, operações com polinômios).
\end{itemize}
} \\ \hline
\multicolumn{3}{|p{\dimexpr\linewidth-2\tabcolsep-2\arrayrulewidth\relax}|}{\textcolor{ifscgreen}{\textbf{Estratégias de Ensino e Aprendizagem:}}} \\ \hline
\multicolumn{3}{|p{\dimexpr\linewidth-2\tabcolsep-2\arrayrulewidth\relax}|}{- **Metodologia:** Aulas expositivas dialogadas; aulas de exercícios; apresentação em linguagem verbal por escrito ou em diálogos na sala de aula; discussões em grupos; estudos dirigidos; leitura e interpretação por meio de datashow; pesquisas em laboratório de informática; seminários; trabalho em laboratório de informática e na biblioteca; trabalhos individuais e em grupos; uso de jogos e objetos de aprendizagem; experimentação no ensino de matemática; produção de vídeos de matemática por estudantes.
- **Avaliação:** Avaliações qualitativas e quantitativas durante o semestre.} \\ \hline
\multicolumn{3}{|p{\dimexpr\linewidth-2\tabcolsep-2\arrayrulewidth\relax}|}{\textcolor{ifscgreen}{\textbf{Bibliografia Básica:}}} \\ \hline
\multicolumn{3}{|p{\dimexpr\linewidth-2\tabcolsep-2\arrayrulewidth\relax}|}{1. BONJORNO, José Roberto; GIOVANNI JÚNIOR, José Ruy; GIOVANNI, José Ruy. \textbf{Matemática fundamental: uma nova abordagem: ensino médio: volume único}. 3. ed. São Paulo: FTD, 2013.
2. DANTE, Luiz Roberto. \textbf{Projeto Múltiplo: matemática: ensino médio, 3º ano}. São Paulo: Ática, 2014.
3. GIOVANNI, José Ruy et al. \textbf{360º matemática fundamental: uma nova abordagem: volume único}. 2. ed. São Paulo: FTD, 2015.} \\ \hline
\multicolumn{3}{|p{\dimexpr\linewidth-2\tabcolsep-2\arrayrulewidth\relax}|}{\textcolor{ifscgreen}{\textbf{Bibliografia Complementar:}}} \\ \hline
\multicolumn{3}{|p{\dimexpr\linewidth-2\tabcolsep-2\arrayrulewidth\relax}|}{1. LIMA, E. L.; CARVALHO, Paulo Cesar P. \textbf{A Matemática do ensino médio}. 7. ed. Rio de Janeiro: SBM, 2016.
2. PAIVA, Manoel. \textbf{Matemática: Paiva: 3º ano}. 2. ed. São Paulo: Moderna, 2013.} \\ \hline
\end{xltabular}

\vspace{0.4cm}

